\documentclass[a4paper,12pt]{article}

% --- Packages essentiels ---
\usepackage[utf8]{inputenc}
\usepackage[T1]{fontenc}
\usepackage[french]{babel}
\usepackage[margin=2.5cm]{geometry}
\usepackage{hyperref}
\usepackage{xcolor}
\usepackage{listings}
\usepackage{tcolorbox}
\usepackage{titlesec}
\usepackage{fancyhdr}
\usepackage{amsmath,amssymb}
\usepackage{graphicx}

% --- Couleurs personnalisées ---
\definecolor{primary}{RGB}{30, 60, 114}
\definecolor{secondary}{RGB}{42, 82, 152}
\definecolor{accent}{RGB}{235, 87, 87}
\definecolor{darkgray}{RGB}{50, 50, 50}
\definecolor{codebg}{RGB}{245, 247, 250}
\definecolor{codeframe}{RGB}{210, 215, 225}

% --- Configuration des liens ---
\hypersetup{
    colorlinks=true,
    linkcolor=primary,
    filecolor=secondary,      
    urlcolor=primary,
    pdftitle={Rapport TP 2 DevOps - Travail collaboratif avec GitLab},
    pdfpagemode=FullScreen,
}

% --- Configuration des listings de code ---
\lstset{
    backgroundcolor=\color{codebg},
    basicstyle=\ttfamily\small,
    breakatwhitespace=false,
    breaklines=true,
    captionpos=b,
    commentstyle=\color{gray},
    frame=single,
    rulecolor=\color{codeframe},
    keepspaces=true,
    keywordstyle=\color{primary}\bfseries,
    language=Python,
    numbers=left,
    numbersep=8pt,
    numberstyle=\tiny\color{gray},
    showspaces=false,
    showstringspaces=false,
    showtabs=false,
    stringstyle=\color{accent},
    tabsize=4
}

% --- En-têtes et pieds de page ---
\pagestyle{fancy}
\fancyhf{}
\rhead{\color{gray}\small TP 2 : Travail Collaboratif avec GitLab}
\lhead{\color{gray}\small DevOps - FastAPI}
\cfoot{\thepage}

\begin{document}

% ==============================================================================
% PAGE DE GARDE
% ==============================================================================
\begin{titlepage}
    \centering
    \vspace*{1cm}
    
    {\Huge \textbf{RAPPORT DE TRAVAIL PRATIQUE}}\\[0.4cm]
    {\LARGE \textbf{TP 2 : Travail Collaboratif avec GitLab}}\\[0.2cm]
    {\large \textbf{Intégration Continue \& Plateforme Collaborative (DevOps)}}\\[1.5cm]
    
    \rule{\linewidth}{0.5mm}\\[0.4cm]
    {\LARGE \textbf{Développement et Fusion d'une API REST FastAPI en Équipe}}\\[0.2cm]
    \rule{\linewidth}{0.5mm}\\[2cm]

    \begin{minipage}{0.48\textwidth}
        \textbf{\large Membres du Groupe :}\\[0.3cm]
        \begin{itemize}
            \item \textbf{Leader :} Mohsen BELLAAJ
            \item \textbf{Dev 1 :} Yessin
            \item \textbf{Dev 2 :} Ahmed
        \end{itemize}
    \end{minipage}
    \hfill
    \begin{minipage}{0.48\textwidth}
        \textbf{\large Cadre Pédagogique :}\\[0.3cm]
        \textbf{Responsable de cours :} Bassem BOUAZIZ\\
        \textbf{Enseignante de TP :} Siwar CHAABENE\\
        \textbf{Auditoires :} D-INFO / D-IITWM\\
        \textbf{Année Universitaire :} 2026-2027
    \end{minipage}

    \vfill
    {\large \today}
\end{titlepage}

\tableofcontents
\newpage

% ==============================================================================
% 1. INTRODUCTION & OBJECTIFS
% ==============================================================================
\section{Introduction et Objectifs du TP}

Le présent travail pratique s'inscrit dans le cadre du module \textbf{Plateforme collaborative et intégration continue (DevOps)}. L'objectif principal est de mettre en application les pratiques modernes d'ingénierie logicielle collaborative en utilisant \textbf{Git} et \textbf{GitLab}, tout en concevant une API REST de référence avec le framework moderne \textbf{FastAPI}.

\subsection{Objectifs pédagogiques et techniques}
\begin{itemize}
    \item \textbf{Architecture logicielle} : Développer un service web RESTful modulaire et léger avec Python et FastAPI.
    \item \textbf{Gestion des versions et branches} : Mettre en œuvre le modèle de branchement Git (Git Feature Branch Workflow) avec la branche principale \texttt{main} et des branches de fonctionnalités isolées.
    \item \textbf{Revue de code et intégration continue} : Utiliser les Merge Requests (MR) sur GitLab pour assurer la validation croisée des ajouts de code.
    \item \textbf{Gestion des conflits} : Détecter, comprendre et résoudre méthodiquement les conflits de fusion survenant lors de développements concurrents sur le même fichier source.
    \item \textbf{Livraison logicielle} : Étiqueter les versions stables via le système de tags Git (\texttt{v0.2}).
\end{itemize}

% ==============================================================================
% 2. ORGANISATION DE L'ÉQUIPE ET RÔLES
% ==============================================================================
\section{Organisation de l'Équipe et Répartition des Rôles}

Conformément aux instructions du TP, la répartition des rôles a été établie sur la base de l'ordre alphabétique des prénoms des membres du groupe :

\begin{table}[h!]
\centering
\begin{tabular}{|l|l|p{8.5cm}|}
\hline
\textbf{Rôle} & \textbf{Étudiant} & \textbf{Responsabilités et Tâches Principales} \\
\hline
\textbf{Leader} & \textbf{Mohsen} & 
\begin{itemize}
    \item Création et configuration du projet GitLab.
    \item Gestion des droits et des invitations des membres.
    \item Initialisation du socle FastAPI sur la branche \texttt{main}.
    \item Revue et approbation des Merge Requests.
    \item Supervision de la résolution de conflits et pose du tag de version.
\end{itemize} \\
\hline
\textbf{Dev 1} & \textbf{Yessin} & 
\begin{itemize}
    \item Clonage et configuration de l'environnement de développement.
    \item Création de la branche \texttt{feature/addition}.
    \item Implémentation du endpoint \texttt{GET /addition}.
    \item Soumission de la Merge Request 1 vers \texttt{main}.
\end{itemize} \\
\hline
\textbf{Dev 2} & \textbf{Ahmed} & 
\begin{itemize}
    \item Clonage et configuration de l'environnement de développement.
    \item Création de la branche \texttt{feature/greeting}.
    \item Implémentation du endpoint \texttt{GET /greeting} (avec Hostname \& IP).
    \item Soumission de la Merge Request 2 et collaboration sur le conflit.
\end{itemize} \\
\hline
\end{tabular}
\caption{Répartition des responsabilités au sein de l'équipe}
\end{table}

% ==============================================================================
% 3. DÉROULEMENT DÉTAILLÉ DES ÉTAPES DU TP
% ==============================================================================
\section{Déroulement Détaillé des Étapes Réalisées}

\subsection{Étape 1 : Création du projet GitLab et Attribution des Rôles}
\textbf{Intervenant :} Leader (Mohsen) avec l'assistance de Dev 1 (Yessin) et Dev 2 (Ahmed).

\begin{enumerate}
    \item Création d'un nouveau projet public/privé sur GitLab intitulé \texttt{tp2-fastapi-team} (ou \texttt{TP2-DEV}).
    \item Accès au menu \textit{Manage $\rightarrow$ Members $\rightarrow$ Invite members}.
    \item Ajout de \textbf{Yessin} (Dev 1) et \textbf{Ahmed} (Dev 2) avec le rôle \textbf{Developer}, leur conférant les droits de création de branches et de push sur leurs branches dédiées.
\end{enumerate}

\subsection{Étape 2 : Vérification et Validation des Accès}
\textbf{Intervenants :} Dev 1 (Yessin) et Dev 2 (Ahmed).

Les développeurs ont vérifié la réception de l'e-mail de confirmation d'invitation envoyé par GitLab et ont validé leur accès au tableau de bord du projet avec les privilèges \textit{Developer}.

\subsection{Étape 3 : Clonage et Configuration de l'Environnement de Développement}
\textbf{Intervenants :} Tous les membres de l'équipe.

\begin{enumerate}
    \item \textbf{Génération du Personal Access Token (PAT)} : Chaque développeur a configuré son jeton d'accès sécurisé dans les préférences GitLab (\textit{Settings $\rightarrow$ Access Tokens}) avec les permissions \texttt{read\_repository} et \texttt{write\_repository}.
    \item \textbf{Clonage du dépôt} :
\begin{lstlisting}[language=bash]
git clone https://gitlab.com/bellajmohsen7/tp2-fastapi-team.git
cd tp2-fastapi-team
\end{lstlisting}
    \item \textbf{Création de l'environnement virtuel local} :
\begin{lstlisting}[language=bash]
python -m venv .venv
# Activation sous Windows PowerShell :
.venv\Scripts\activate
\end{lstlisting}
    \item \textbf{Installation des dépendances fondamentales} :
\begin{lstlisting}[language=bash]
pip install fastapi uvicorn
\end{lstlisting}
    \item \textbf{Création du fichier \texttt{.gitignore}} pour exclure l'environnement virtuel et les caches :
\begin{lstlisting}[language=bash]
.venv/
__pycache__/
*.pyc
\end{lstlisting}
\end{enumerate}

\subsection{Étape 4 : Développement du Socle Applicatif sur \texttt{main} (Commit A)}
\textbf{Intervenant :} Leader (Mohsen).

Le leader a implémenté la structure initiale de l'application dans le fichier \texttt{main.py} :
\begin{lstlisting}[language=Python]
from fastapi import FastAPI

app = FastAPI(title="TP2 FastAPI App")

@app.get("/")
def read_root():
    return {"message": "it works !"}
\end{lstlisting}

Il a également rédigé le fichier \texttt{README.md} détaillant :
\begin{itemize}
    \item La procédure d'installation et d'activation de l'environnement.
    \item La commande d'exécution du serveur : \texttt{uvicorn main:app --reload}.
    \item Les modalités de test via navigateur (\url{http://localhost:8000/}) et via la documentation Swagger (\url{http://localhost:8000/docs}).
\end{itemize}

Enfin, le leader a validé et poussé ces ajouts sur la branche centrale :
\begin{lstlisting}[language=bash]
git add .
git commit -m "feat: implement GET / endpoint and update README"
git push origin main
\end{lstlisting}

\subsection{Étape 5 : Récupération et Validation par les Développeurs}
\textbf{Intervenants :} Dev 1 (Yessin) et Dev 2 (Ahmed).

Les deux développeurs ont synchronisé leur copie locale avec la branche distante :
\begin{lstlisting}[language=bash]
git pull origin main
uvicorn main:app --reload
\end{lstlisting}
Le test local sur \url{http://localhost:8000/} a confirmé le bon fonctionnement de l'endpoint racine renvoyant \texttt{\{"message": "it works !"\}}.

\subsection{Étape 6 : Travail Parallèle sur les Branches de Fonctionnalités}

\subsubsection{Travail de Dev 1 (Yessin) : Branche \texttt{feature/addition}}
Yessin a créé sa branche isolée, développé la route d'addition et partagé son travail :
\begin{enumerate}
    \item Création de la branche :
\begin{lstlisting}[language=bash]
git checkout -b feature/addition
\end{lstlisting}
    \item Ajout du code dans \texttt{main.py} :
\begin{lstlisting}[language=Python]
@app.get("/addition")
def addition(a: int, b: int):
    return {
        "result a+b": a + b,
        "developer": "Yessin"
    }
\end{lstlisting}
    \item Validation locale sur \url{http://localhost:8000/addition?a=3&b=7} renvoyant :
\begin{lstlisting}[language=json]
{
  "result a+b": 10,
  "developer": "Yessin"
}
\end{lstlisting}
    \item Commit, push et création de la Merge Request (MR 1) vers \texttt{main} :
\begin{lstlisting}[language=bash]
git add main.py
git commit -m "add addition endpoint"
git push origin feature/addition
\end{lstlisting}
\end{enumerate}

\subsubsection{Travail de Dev 2 (Ahmed) : Branche \texttt{feature/greeting}}
Simultanément, Ahmed est parti de la branche \texttt{main} pour créer sa fonctionnalité :
\begin{enumerate}
    \item Création de la branche :
\begin{lstlisting}[language=bash]
git checkout main
git checkout -b feature/greeting
\end{lstlisting}
    \item Implémentation dans \texttt{main.py} de la route \texttt{/greeting} avec capture dynamique du nom d'hôte et de l'adresse IP via le module standard \texttt{socket} :
\begin{lstlisting}[language=Python]
import socket

def get_ip_address():
    s = socket.socket(socket.AF_INET, socket.SOCK_DGRAM)
    try:
        s.connect(("8.8.8.8", 80))
        ip = s.getsockname()[0]
    except Exception:
        ip = "127.0.0.1"
    finally:
        s.close()
    return ip

@app.get("/greeting")
def greeting(name: str):
    return {
        "message": f"Hello, {name}",
        "developer": "Ahmed",
        "hostname": socket.gethostname(),
        "IP": get_ip_address()
    }
\end{lstlisting}
    \item Validation locale sur \url{http://localhost:8000/greeting?name=Prof}.
    \item Commit, push et ouverture de la Merge Request (MR 2) vers \texttt{main} :
\begin{lstlisting}[language=bash]
git add main.py
git commit -m "add greeting endpoint"
git push origin feature/greeting
\end{lstlisting}
\end{enumerate}

\subsection{Étape 7 : Revue et Fusion de la MR de Dev 1 (Merge 1)}
\textbf{Intervenant :} Leader (Mohsen).

\begin{enumerate}
    \item Ouverture de la Merge Request 1 (\texttt{feature/addition} $\rightarrow$ \texttt{main}).
    \item Revue du code dans l'onglet \textit{Changes}.
    \item Approbation (\textit{Approve}) et exécution de la fusion (\textit{Merge}) avec suppression automatique de la branche distante.
\end{enumerate}

\begin{tcolorbox}[colback=blue!5!white,colframe=primary,title=\textbf{Réponse à la Question : Interpréter le Merge 1}]
La fusion de la branche \texttt{feature/addition} s'est effectuée de façon fluide et \textbf{sans aucun conflit}.\\
\textbf{Explication :} Depuis le moment où Yessin (Dev 1) a tiré sa branche à partir de \texttt{main}, aucun commit supplémentaire n'a été inséré sur la branche principale \texttt{main}. L'historique Git était donc direct et linéaire. Git a pu appliquer la modification sans ambiguïté (Fast-Forward ou Merge commit automatique).
\end{tcolorbox}

\subsection{Étape 8 : Revue de la MR de Dev 2 et Résolution du Conflit (Merge 2)}
\textbf{Intervenants :} Leader (Mohsen), Dev 2 (Ahmed) et Dev 1 (Yessin).

À l'ouverture de la MR 2 soumise par Ahmed, GitLab affiche un blocage majeur :
\begin{center}
\textbf{\color{accent}\fbox{Merge blocked: Merge conflicts must be resolved}}
\end{center}

\begin{tcolorbox}[colback=red!5!white,colframe=accent,title=\textbf{Réponse à la Question : Interpréter le Conflit}]
\textbf{Explication du conflit :} Yessin (Dev 1) et Ahmed (Dev 2) ont tous les deux dérivé leurs branches respectives de la même version originale de \texttt{main}.\\
Entre-temps, la MR 1 de Dev 1 a été fusionnée dans \texttt{main}, introduisant de nouvelles lignes dans \texttt{main.py}.\\
Lorsque Dev 2 tente de fusionner sa branche, Git constate que le même fichier (\texttt{main.py}) a subi des modifications divergentes aux mêmes emplacements de lignes. Ne pouvant deviner s'il faut conserver l'addition, le greeting ou les deux, Git suspend l'intégration et délègue l'arbitrage à l'équipe.
\end{tcolorbox}

\subsubsection{Procédure de Résolution du Conflit par l'Équipe}
La résolution a été opérée en local avec rigueur :
\begin{enumerate}
    \item Récupération des dernières modifications de \texttt{main} sur la branche \texttt{feature/greeting} :
\begin{lstlisting}[language=bash]
git checkout feature/greeting
git pull origin main
\end{lstlisting}
    \item Analyse des délimiteurs de conflit dans \texttt{main.py} :
\begin{lstlisting}[language=Python]
<<<<<<< HEAD (Version d'Ahmed - greeting)
@app.get("/greeting")
def greeting(name: str):
...
=======
@app.get("/addition")
def addition(a: int, b: int):
...
>>>>>>> main (Version de Yessin - addition)
\end{lstlisting}
    \item \textbf{Arbitrage} : Nettoyage manuel des marqueurs pour intégrer \textbf{l'ensemble des deux fonctionnalités} dans le fichier final consolidé :
\begin{lstlisting}[language=Python]
import socket
from fastapi import FastAPI

app = FastAPI(title="TP2 FastAPI App")

def get_ip_address():
    s = socket.socket(socket.AF_INET, socket.SOCK_DGRAM)
    try:
        s.connect(("8.8.8.8", 80))
        ip = s.getsockname()[0]
    except Exception:
        ip = "127.0.0.1"
    finally:
        s.close()
    return ip

@app.get("/")
def read_root():
    return {"message": "it works !"}

@app.get("/addition")
def addition(a: int, b: int):
    return {
        "result a+b": a + b,
        "developer": "Yessin"
    }

@app.get("/greeting")
def greeting(name: str):
    return {
        "message": f"Hello, {name}",
        "developer": "Ahmed",
        "hostname": socket.gethostname(),
        "IP": get_ip_address()
    }
\end{lstlisting}
    \item Validation et publication du correctif de fusion :
\begin{lstlisting}[language=bash]
git add main.py
git commit -m "fix: resolve merge conflict between addition and greeting"
git push origin feature/greeting
\end{lstlisting}
    \item Le statut sur GitLab bascule instantanément au vert. Mohsen (Leader) valide le merge final sur la branche \texttt{main}.
\end{enumerate}

\subsection{Étape Finale : Étiquetage de Version (Tag v0.2)}
\textbf{Intervenant :} Leader (Mohsen).

Après validation complète des trois routes de l'API en production locale, la version de release a été officialisée par la création d'un tag Git :
\begin{lstlisting}[language=bash]
git checkout main
git pull origin main
git tag -a v0.2 -m "Release version 0.2: Consolidated API with addition and greeting endpoints"
git push origin v0.2
\end{lstlisting}

% ==============================================================================
% 4. BILAN TECHNIQUE ET SYNTHÈSE
% ==============================================================================
\section{Bilan Technique et Bonnes Pratiques DevOps}

Ce travail pratique a permis de mettre en lumière des concepts clés de la collaboration logicielle moderne :
\begin{enumerate}
    \item \textbf{Isolation des développements} : L'utilisation systématique de branches dédiées (\texttt{feature/*}) préserve la stabilité de la branche principale \texttt{main}.
    \item \textbf{Revue par les pairs (Code Review)} : Les Merge Requests ne sont pas de simples outils de fusion, mais des garde-fous garantissant la qualité du code avant intégration.
    \item \textbf{Gestion saine des conflits} : Les conflits Git font partie intégrante de la vie d'un projet d'équipe ; leur résolution méthodique en équipe renforce la cohésion et la compréhension globale de la base de code.
    \item \textbf{Traçabilité par les tags} : L'étiquetage sémantique (\texttt{v0.2}) pose les fondations des déclencheurs de déploiement continu (CD) qui seront explorés dans les prochains TPs.
\end{enumerate}

\end{document}
